
\documentclass[11pt]{article}
\usepackage[margin=1in]{geometry}
\usepackage{amsmath,amssymb,amsthm,mathtools}
\usepackage[T1]{fontenc}
\usepackage[utf8]{inputenc}
\usepackage{enumitem}
\usepackage{booktabs,longtable,array}
\usepackage{hyperref}
\usepackage{xcolor}
\usepackage{microtype}
\hypersetup{colorlinks=true, linkcolor=blue!50!black, citecolor=blue!50!black, urlcolor=blue!50!black}

\newtheorem{theorem}{Theorem}[section]
\newtheorem{proposition}[theorem]{Proposition}
\newtheorem{lemma}[theorem]{Lemma}
\newtheorem{corollary}[theorem]{Corollary}
\newtheorem{claim}[theorem]{Claim}
\theoremstyle{definition}
\newtheorem{definition}[theorem]{Definition}
\newtheorem{assumption}[theorem]{Assumption}
\theoremstyle{remark}
\newtheorem{remark}[theorem]{Remark}

\newcommand{\KL}{\mathrm{KL}}
\newcommand{\Ent}{\mathrm{Ent}}
\newcommand{\Prob}{\mathsf{P}}
\newcommand{\E}{\mathbb{E}}
\newcommand{\R}{\mathbb{R}}
\newcommand{\N}{\mathbb{N}}
\newcommand{\Pcal}{\mathcal{P}}
\newcommand{\Mcal}{\mathcal{M}}
\newcommand{\Qcal}{\mathcal{Q}}
\newcommand{\Zcal}{\mathcal{Z}}
\newcommand{\Fcal}{\mathcal{F}}
\newcommand{\Gcal}{\mathcal{G}}
\newcommand{\W}{\mathcal{W}}
\newcommand{\one}{\mathbf{1}}
\newcommand{\dd}{\,d}
\newcommand{\supp}{\operatorname{supp}}
\newcommand{\argmin}{\operatorname*{arg\,min}}
\newcommand{\argmax}{\operatorname*{arg\,max}}
\newcommand{\esssup}{\operatorname*{ess\,sup}}
\newcommand{\ri}{\operatorname{ri}}
\newcommand{\dom}{\operatorname{dom}}
\newcommand{\Law}{\operatorname{Law}}
\newcommand{\Var}{\operatorname{Var}}
\newcommand{\Cov}{\operatorname{Cov}}
\newcommand{\diag}{\operatorname{diag}}
\newcommand{\Id}{\mathrm{Id}}
\newcommand{\Lip}{\operatorname{Lip}}
\newcommand{\osc}{\operatorname{osc}}
\newcommand{\diam}{\operatorname{diam}}
\newcommand{\proj}{\operatorname{proj}}
\newcommand{\e}{\mathrm{e}}

\newenvironment{argument}{\paragraph{Argument.}\begin{enumerate}[leftmargin=2em]}{\end{enumerate}}
\newenvironment{proofoutline}{\paragraph{Proof architecture.}\begin{enumerate}[leftmargin=2em]}{\end{enumerate}}


% Source-faithfulness apparatus used by the paper reconstructions.
\newcommand{\sourceanchor}[1]{\par\smallskip\noindent\textbf{Source anchor.} #1\par\smallskip}
\newcommand{\sourcecomment}[1]{\begin{remark}#1\end{remark}}
\newenvironment{sourceguide}
  {\begin{description}[style=nextline,leftmargin=3.2cm,labelwidth=3cm]}
  {\end{description}}

% Backward-compatible environments used by some of the older reconstructions.
\newenvironment{distilled}[1][]{\begin{theorem}[#1]}{\end{theorem}}
\newenvironment{distillation}{\begin{enumerate}[leftmargin=2em]}{\end{enumerate}}

\newenvironment{commentarynotes}{\begin{remark}\begin{enumerate}[leftmargin=2em]}{\end{enumerate}\end{remark}}

\title{Step-by-step modernization of Franklin--Lorenz (1989), Section 3: Hilbert metric proof of Sinkhorn iteration convergence}
\author{}
\date{Source PDF: \texttt{FranklinLorenz1989-1-s2.0-0024379589904904-main.pdf}}
\begin{document}
\maketitle

\section*{Source map and scope}
\begin{description}[style=nextline,leftmargin=3.2cm,labelwidth=3cm]
\item[Primary source] \texttt{\detokenize{FranklinLorenz1989-1-s2.0-0024379589904904-main.pdf}}
\item[Coverage] Section 3 in full: Sinkhorn's alternating normalization, Hilbert's projective metric, Lemmas 1--3, Theorem 4, equations (3.9) and (3.12), and both numerical examples.
\item[Source anchors]
\begin{itemize}[leftmargin=2em]
  \item Section 3
  \item Lemma 1
  \item Lemma 2
  \item Theorem 4
  \item Lemma 3
  \item equations (3.9) and (3.12)
  \item Examples 1--2 and Tables 1--2
\end{itemize}
\end{description}

The exposition below preserves the source proof order and proof technique.  Notation is normalized
and intermediate calculations are made explicit.  When the source contains a gap, ambiguity, or
apparently incorrect step, the reconstruction records the source step before adding a separate
commentary remark.

\section*{Reference and scope}
This note modernizes the final section of Joel Franklin and Jens Lorenz,
\emph{On the Scaling of Multidimensional Matrices}, Linear Algebra and its
Applications 114/115:717--735, 1989 \cite{FranklinLorenz1989}.  The portion
expanded here is Section 3, ``Hilbert's Projective Metric and Sinkhorn's
Iteration'', pp. 725--732 of the article.  Earlier sections prove existence of
scaled multidimensional matrices by convex minimization.  Section 3 is narrower:
it revisits Sinkhorn's original two-dimensional positive-matrix iteration and gives
a quantitative geometric convergence estimate from Hilbert's projective metric and
Birkhoff's contraction theorem.

The goal of this writeup is not to compress the argument to its slogan.  It keeps
the algebraic identities that make the proof work.  The key load-bearing pieces are:
\begin{enumerate}[leftmargin=2em]
\item the row/column normalization definitions;
\item the Hilbert metric and Birkhoff contraction input;
\item Lemma 2, which expresses one normalization cycle as two positive linear maps;
\item Theorem 4, which turns Lemma 2 into a geometric Cauchy-tail bound in the
      diagonal-equivalence metric $\Delta$;
\item Lemma 3, which translates $\Delta$-smallness into entrywise multiplicative
      error bounds.
\end{enumerate}

\section*{Section 3 notation}
Let $A=(a_{ij})$ be a strictly positive $m\times n$ matrix.  Let
$p\in\R^m_{>0}$ and $q\in\R^n_{>0}$ have the same total mass,
\[
  \sum_{i=1}^m p_i = \sum_{j=1}^n q_j.
\]
Write $e$ for the all-ones vector, with dimension inferred from context.  The
iteration stays inside the positive diagonal-equivalence class of $A$.

Starting from $A_0=A$, define two alternating sequences:
\[
  r^{(k)} := A_k e \in \R^m_{>0},
  \qquad
  S_k := \operatorname{diag}\left(\frac{p_i}{r_i^{(k)}}\right),
  \qquad
  A'_k := S_k A_k,
\]
so $A'_k$ is row-normalized:
\[
  A'_k e = S_k A_k e = S_k r^{(k)} = p.
\]
Then define
\[
  c^{(k)} := (A'_k)^T e \in \R^n_{>0},
  \qquad
  T_k := \operatorname{diag}\left(\frac{q_j}{c_j^{(k)}}\right),
  \qquad
  A_{k+1} := A'_k T_k,
\]
so $A_{k+1}$ is column-normalized:
\[
  A_{k+1}^T e = T_k (A'_k)^T e = T_k c^{(k)} = q.
\]
Franklin--Lorenz's Theorem 4 assumes the initial matrix is already column-normalized:
\[
  A^T e=q.
\]
Consequently $A_k^T e=q$ for every $k\ge 0$ after the first column-normalized input,
and $A'_k e=p$ for every $k\ge 0$ by construction.

\section*{Hilbert projective metric and the Birkhoff input}
For positive vectors $x,x'\in\R^m_{>0}$ Franklin--Lorenz use Hilbert's projective
metric
\[
  d(x,x')=
  \log \max_{i,k}\frac{x_i x'_k}{x'_i x_k}.
\]
Equivalently,
\[
  d(x,x')=
  \log\frac{\max_i(x_i/x'_i)}{\min_i(x_i/x'_i)}.
\]
This has the usual projective properties:
\begin{align*}
  d(x,x')=0 &\iff x=\alpha x'\text{ for some }\alpha>0,\\
  d(x,x') &= d(x',x),\\
  d(x,x'') &\le d(x,x')+d(x',x'').
\end{align*}
It is invariant under independent positive rescaling of either vector:
\[
  d(\alpha x,\beta x')=d(x,x')\qquad(\alpha,\beta>0).
\]
The simple quotient identity used repeatedly in the paper is
\[
  d(x/x',e)=d(x,x'),
\]
where $x/x'$ denotes coordinatewise quotient.  Indeed,
\[
  d(x/x',e)=\log\frac{\max_i(x_i/x'_i)}{\min_i(x_i/x'_i)}=d(x,x').
\]

For a positive matrix $A$, Franklin--Lorenz write a projective diameter parameter
$\vartheta(A)$.  Their Lemma 1 is the exact Birkhoff contraction formula:
\begin{lemma}[Franklin--Lorenz Lemma 1, source pp.~727--728]
For a strictly positive matrix $A$, let
\[
  \kappa(A)=
  \sup_{y\not\sim y'}\frac{d(Ay,Ay')}{d(y,y')}.
\]
Then
\[
  \kappa(A)=
  \frac{\sqrt{\vartheta(A)}-1}{\sqrt{\vartheta(A)}+1}.
\]
\end{lemma}

With the logarithmic convention above, the clean modern reading is
that $\vartheta(A)$ is the multiplicative Hilbert diameter:
\[
  \vartheta(A)
  := \max_{i,j,k,\ell}\frac{a_{ik}a_{j\ell}}{a_{jk}a_{i\ell}}
  = \exp\left(\sup_{y,y'>0} d(Ay,Ay')\right).
\]
The printed paper writes this immediately next to the additive metric $d$; the
formula used in Lemma 1 is the multiplicative one.  The paper attributes the result
to Birkhoff and also points to Bauer and Bushell
\cite{Birkhoff1957,Bauer1965,Bushell1973}.  The contraction coefficient is
\[
  \kappa(A)
  =\sup_{y\not\sim y'}\frac{d(Ay,Ay')}{d(y,y')}
  =\frac{\vartheta(A)^{1/2}-1}{\vartheta(A)^{1/2}+1}.
\]
Strict positivity and finite dimension give $1\le \vartheta(A)<\infty$, hence
$0\le \kappa(A)<1$ unless the image is projectively a point.  The estimates below
only need that $\kappa(A)<1$.

Two invariance facts are used without ceremony in the paper.  If $B=XAY$ with
$X,Y$ positive diagonal, then every cross-ratio in the entries is unchanged:
\[
  \frac{(x_i a_{ik}y_k)(x_j a_{j\ell}y_\ell)}
       {(x_j a_{jk}y_k)(x_i a_{i\ell}y_\ell)}
  =\frac{a_{ik}a_{j\ell}}{a_{jk}a_{i\ell}}.
\]
Therefore
\[
  A\sim B \quad\Longrightarrow\quad
  \vartheta(A)=\vartheta(B),\qquad \kappa(A)=\kappa(B).
\]
Also, transposition leaves the same set of cross-ratios, so
\[
  \vartheta(A^T)=\vartheta(A),\qquad \kappa(A^T)=\kappa(A).
\]
Finally, composition cannot increase contraction constants beyond the product:
\[
  \kappa(A_1A_2)\le \kappa(A_1)\kappa(A_2),
\]
whenever the product is defined and the matrices are positive.

\section*{Lemma 2: one Sinkhorn cycle contracts marginal errors}
Franklin--Lorenz's Lemma 2 is the central algebraic observation.  It states:

\begin{quote}
Assume $A=A_0$ is column-normalized, $A^T e=q$.  Let $A'_0=S_0A_0$ and
$A_1=A'_0T_0$ be the first row and column normalization steps.  Then
\[
  d(r^{(1)},p)\le \gamma d(r^{(0)},p),
  \qquad
  d(c^{(1)},q)\le \gamma d(c^{(0)},q),
\]
where
\[
  \gamma:=\kappa(A)^2.
\]
\end{quote}

The proof is short in the paper, but every identity matters.

\subsection*{First half: row-error estimate}
Set $A'=A'_0=S_0A$.  Since $A_1=A'T_0$ and $T_0e=q/c^{(0)}$, the next row-sum vector is
\[
  r^{(1)}=A_1e=A'T_0e=A'\left(\frac{q}{c^{(0)}}\right).
\]
Since $A'$ is row-normalized,
\[
  A'e=p.
\]
Thus the row error after the full cycle is
\[
  d(r^{(1)},p)
  =d\left(A'\frac{q}{c^{(0)}},A'e\right).
\]
Apply Birkhoff's contraction theorem to the positive matrix $A'$:
\[
  d\left(A'\frac{q}{c^{(0)}},A'e\right)
  \le \kappa(A')d\left(\frac{q}{c^{(0)}},e\right).
\]
Use the quotient identity:
\[
  d\left(\frac{q}{c^{(0)}},e\right)=d(q,c^{(0)}).
\]
Because $A'$ is diagonally equivalent to $A$, $\kappa(A')=\kappa(A)$.  Hence
\[
  d(r^{(1)},p)\le \kappa(A)d(q,c^{(0)}).
\]
It remains to bound $d(q,c^{(0)})$ by another factor of $\kappa(A)$ times the old row
error.  Compute the column sums of $A'=S_0A$:
\[
  c^{(0)}=(A')^Te=A^T S_0e.
\]
But $S_0e=p/r^{(0)}$, so
\[
  c^{(0)}=A^T\left(\frac{p}{r^{(0)}}\right).
\]
The initial column-normalization hypothesis is $A^Te=q$, hence
\[
  d(q,c^{(0)})
  =d\left(A^Te,A^T\frac{p}{r^{(0)}}\right)
  \le \kappa(A^T)d\left(e,\frac{p}{r^{(0)}}\right).
\]
Again use transposition invariance $\kappa(A^T)=\kappa(A)$ and the quotient identity:
\[
  d\left(e,\frac{p}{r^{(0)}}\right)=d\left(\frac{p}{r^{(0)}},e\right)=d(p,r^{(0)})=d(r^{(0)},p).
\]
Combining the two displayed inequalities gives
\[
  d(r^{(1)},p)
  \le \kappa(A)^2 d(r^{(0)},p)
  =\gamma d(r^{(0)},p).
\]

\subsection*{Second half: column-error estimate}
The paper says that the column estimate follows similarly.  Here is the expanded
version.

After the first full cycle, $A_1$ is column-normalized:
\[
  A_1^Te=q.
\]
The next row-normalized matrix is
\[
  A'_1=S_1A_1,
  \qquad
  S_1e=\frac{p}{r^{(1)}}.
\]
Therefore its column sums are
\[
  c^{(1)}=(A'_1)^Te=A_1^T S_1e
  =A_1^T\left(\frac{p}{r^{(1)}}\right).
\]
Since $A_1^Te=q$, Birkhoff contraction for $A_1^T$ gives
\[
  d(c^{(1)},q)
  =d\left(A_1^T\frac{p}{r^{(1)}},A_1^Te\right)
  \le \kappa(A_1^T)d\left(\frac{p}{r^{(1)}},e\right).
\]
Use diagonal-equivalence and transposition invariance:
\[
  \kappa(A_1^T)=\kappa(A_1)=\kappa(A).
\]
Use the quotient identity:
\[
  d\left(\frac{p}{r^{(1)}},e\right)=d(p,r^{(1)})=d(r^{(1)},p).
\]
Thus
\[
  d(c^{(1)},q)\le \kappa(A)d(r^{(1)},p).
\]
But the first part of the proof, stopping before the second application through
$A^T$, already gave the sharper relation
\[
  d(r^{(1)},p)\le \kappa(A)d(c^{(0)},q).
\]
Therefore
\[
  d(c^{(1)},q)
  \le \kappa(A)^2 d(c^{(0)},q)
  =\gamma d(c^{(0)},q).
\]
This completes Lemma 2.

\section*{Iterating Lemma 2}
Every $A_k$ with $k\ge 1$ is column-normalized by construction.  Theorem 4 assumes
$A_0$ is also column-normalized.  Therefore Lemma 2 can be restarted at each $A_k$.
For all $k\ge 0$,
\[
  d(r^{(k+1)},p)\le \gamma d(r^{(k)},p),
  \qquad
  d(c^{(k+1)},q)\le \gamma d(c^{(k)},q).
\]
By induction,
\[
  d(r^{(k)},p)\le \gamma^k d(r^{(0)},p),
  \qquad
  d(c^{(k)},q)\le \gamma^k d(c^{(0)},q).
\]
The constant $\gamma$ is strictly less than $1$ because $A$ is strictly positive and
finite-dimensional, so $\kappa(A)<1$.

\section*{The metric $\Delta$ on the diagonal-equivalence class}
Let
\[
  E_A:=\{B:B\sim A,\; b_{ij}>0\}
\]
be the positive diagonal-equivalence class of $A$.  If $B,B'\in E_A$ and
\[
  B=XB'Y,
  \qquad
  X=\operatorname{diag}(x_i),
  \qquad
  Y=\operatorname{diag}(y_j),
\]
Franklin--Lorenz define
\[
  \Delta(B,B'):=d(x,e)+d(y,e).
\]
The diagonal factors are not unique: replacing $X$ by $\lambda X$ and $Y$ by
$\lambda^{-1}Y$ gives the same product.  But $d(\lambda x,e)=d(x,e)$ and
$d(\lambda^{-1}y,e)=d(y,e)$, so $\Delta$ is independent of this scalar ambiguity.
The paper states that this is a metric on $E_A$ and that $(E_A,\Delta)$ is complete.
For the proof of Theorem 4 as printed, only the triangle inequality and the already
known convergence of Sinkhorn's iteration are used.

For consecutive Sinkhorn iterates,
\[
  A_{k+1}=S_kA_kT_k,
  \qquad
  S_ke=\frac{p}{r^{(k)}},
  \qquad
  T_ke=\frac{q}{c^{(k)}}.
\]
Therefore
\begin{align*}
  \Delta(A_k,A_{k+1})
  &=d\left(\frac{p}{r^{(k)}},e\right)
    +d\left(\frac{q}{c^{(k)}},e\right)\\
  &=d(r^{(k)},p)+d(c^{(k)},q).
\end{align*}
This exact increment identity is the bridge from marginal errors to matrix convergence.

\section*{Theorem 4: geometric convergence in $\Delta$}
Let $\widehat B$ be the unique positive matrix diagonally equivalent to $A$ with
row sums $p$ and column sums $q$:
\[
  \widehat B e=p,
  \qquad
  \widehat B^T e=q,
  \qquad
  \widehat B\sim A.
\]
Franklin--Lorenz's Theorem 4 states that if $A^Te=q$, then
\[
  \Delta(A_k,\widehat B)
  \le \frac{1}{1-\gamma}
        \{d(r^{(k)},p)+d(c^{(k)},q)\}
  \le \frac{\gamma^k}{1-\gamma}
        \{d(r^{(0)},p)+d(c^{(0)},q)\},
\]
and similarly
\[
  \Delta(A'_k,\widehat B)
  \le \frac{1}{1-\gamma}
        \{d(r^{(k+1)},p)+d(c^{(k)},q)\}
  \le \frac{\gamma^k}{1-\gamma}
        \{d(r^{(1)},p)+d(c^{(0)},q)\}.
\]

\subsection*{Proof for $A_k$}
From the increment identity and the iterated Lemma 2 estimate,
\begin{align*}
  \Delta(A_k,A_{k+1})
  &=d(r^{(k)},p)+d(c^{(k)},q)\\
  &\le \gamma^k\{d(r^{(0)},p)+d(c^{(0)},q)\}.
\end{align*}
More generally, restarting the same estimate at index $k$ gives, for $\ell\ge k$,
\[
  d(r^{(\ell)},p)+d(c^{(\ell)},q)
  \le \gamma^{\ell-k}\{d(r^{(k)},p)+d(c^{(k)},q)\}.
\]
By the triangle inequality,
\[
  \Delta(A_k,A_N)
  \le \sum_{\ell=k}^{N-1}\Delta(A_\ell,A_{\ell+1}).
\]
Letting $N\to\infty$ and using the known convergence $A_N\to\widehat B$ in
$E_A$ gives
\begin{align*}
  \Delta(A_k,\widehat B)
  &\le \sum_{\ell=k}^{\infty}\Delta(A_\ell,A_{\ell+1})\\
  &\le \sum_{j=0}^{\infty}\gamma^j
       \{d(r^{(k)},p)+d(c^{(k)},q)\}\\
  &=\frac{1}{1-\gamma}\{d(r^{(k)},p)+d(c^{(k)},q)\}.
\end{align*}
Applying the global estimates
$d(r^{(k)},p)\le \gamma^kd(r^{(0)},p)$ and
$d(c^{(k)},q)\le \gamma^kd(c^{(0)},q)$ gives the second displayed bound.

\subsection*{Proof for $A'_k$}
The row-normalized matrix $A'_k$ differs from $A_{k+1}$ only by the column correction
$T_k$:
\[
  A_{k+1}=A'_kT_k,
  \qquad
  T_ke=\frac{q}{c^{(k)}}.
\]
Thus
\[
  \Delta(A'_k,A_{k+1})=d\left(\frac{q}{c^{(k)}},e\right)=d(c^{(k)},q).
\]
Then
\[
  \Delta(A'_k,\widehat B)
  \le d(c^{(k)},q)+\Delta(A_{k+1},\widehat B).
\]
Using the already proved $A_k$ bound from index $k+1$,
\[
  \Delta(A_{k+1},\widehat B)
  \le \frac{1}{1-\gamma}
      \{d(r^{(k+1)},p)+d(c^{(k+1)},q)\}.
\]
Since $d(c^{(k+1)},q)\le \gamma d(c^{(k)},q)$,
\begin{align*}
  \Delta(A'_k,\widehat B)
  &\le d(c^{(k)},q)
      +\frac{d(r^{(k+1)},p)+d(c^{(k+1)},q)}{1-\gamma}\\
  &\le \frac{d(r^{(k+1)},p)+d(c^{(k)},q)}{1-\gamma}.
\end{align*}
Finally,
\[
  d(r^{(k+1)},p)\le \gamma^k d(r^{(1)},p),
  \qquad
  d(c^{(k)},q)\le \gamma^k d(c^{(0)},q),
\]
which gives Franklin--Lorenz's second estimate for $A'_k$.

\section*{Completeness note}
The paper explicitly notes that the printed proof uses Sinkhorn's earlier convergence
theorem to identify the limiting matrix $\widehat B$.  But the same estimates also give
an independent route: if one proves that $(E_A,\Delta)$ is complete, then the finite-tail
bound
\[
  \Delta(A_k,A_N)\le
  \frac{\gamma^k}{1-\gamma}
  \{d(r^{(0)},p)+d(c^{(0)},q)\}
\]
shows that $(A_k)$ is Cauchy in $\Delta$, hence converges to some $B\in E_A$; the
marginal-error convergence then forces $Be=p$ and $B^Te=q$, so $B=\widehat B$ by
uniqueness of the scaled matrix.

\section*{Lemma 3: entrywise multiplicative error from $\Delta$}
Franklin--Lorenz's Lemma 3 is an error-bound conversion lemma.  It states:

\begin{quote}
Let $A\sim B$ and assume $A^Te=B^Te$.  If $\Delta(A,B)\le \varepsilon$, then for all
entries
\[
  \exp(-\varepsilon)\le \frac{b_{ij}}{a_{ij}}\le \exp(\varepsilon).
\]
\end{quote}

Here is the full proof.
Since $A\sim B$, choose positive diagonal matrices
\[
  X=\operatorname{diag}(x_i),\qquad Y=\operatorname{diag}(y_j)
\]
such that
\[
  B=XAY.
\]
The condition $\Delta(A,B)\le \varepsilon$ says, after choosing diagonal factors,
\[
  d(x,e)\le \varepsilon,
  \qquad
  d(y,e)\le \varepsilon
\]
up to the harmless scalar ambiguity between $X$ and $Y$.  Franklin--Lorenz scale
$X$ so that its smallest diagonal entry is $1$.  Then
\[
  1\le x_i\le \exp(\varepsilon)
  \qquad(i=1,\dots,m),
\]
and therefore
\[
  \exp(-\varepsilon)e\le X^{-1}e\le e
\]
componentwise.

The column-sum equality $A^Te=B^Te=q$ gives
\[
  Yq=YA^Te=(AY)^Te=(X^{-1}B)^Te=B^TX^{-1}e.
\]
Multiplying the componentwise bound on $X^{-1}e$ by the positive matrix $B^T$ yields
\[
  \exp(-\varepsilon)B^Te
  \le B^TX^{-1}e
  \le B^Te.
\]
Since $B^Te=q$, this is
\[
  \exp(-\varepsilon)q\le B^TX^{-1}e\le q.
\]
Using $Yq=B^TX^{-1}e$ and $q_j>0$ gives
\[
  \exp(-\varepsilon)\le y_j\le 1.
\]
Finally,
\[
  \frac{b_{ij}}{a_{ij}}=x_i y_j.
\]
Combining $1\le x_i\le \exp(\varepsilon)$ and
$\exp(-\varepsilon)\le y_j\le 1$ gives
\[
  \exp(-\varepsilon)\le x_i y_j\le \exp(\varepsilon),
\]
which is the desired entrywise estimate.


\section*{Computable entrywise error bounds: equations (3.9) and (3.12)}
Combining Theorem 4 with Lemma 3 gives a completely computable stopping
criterion.  Define
\[
  \varepsilon_k:=
  \frac{\gamma^k}{1-\gamma}
  \left\{d(r^{(0)},p)+d(c^{(0)},q)\right\},
  \qquad
  \lambda_k:=e^{\varepsilon_k}.
  \tag{3.12}
\]
Then equation (3.9) reads
\[
  \lambda_k^{-1}
  \le \frac{\widehat b_{ij}}{a^{(k)}_{ij}}
  \le \lambda_k
  \qquad\text{for every }i,j.
  \tag{3.9}
\]
The point is that $\lambda_k$ uses only the current data and the a priori
contraction factor.  No knowledge of the unknown scaling matrices is required.
Franklin--Lorenz compare this bound with the actual smallest number
$\widehat\lambda_k$ satisfying the same two-sided inequality.

\section*{Example 1: a slowly contracting matrix}
Take
\[
  p=q=\left(\frac13,\frac13,\frac13\right)^T,
  \qquad
  A=A_0=\frac1{30}
  \begin{pmatrix}
    1&3&8\\
    1&4&1\\
    8&3&1
  \end{pmatrix}.
\]
The limiting doubly stochastic scaling printed in the paper is
\[
  \widehat B=
  \begin{pmatrix}
    0.029629630&0.066666667&0.237037037\\
    0.066666667&0.200000000&0.066666667\\
    0.237037037&0.066666667&0.029629630
  \end{pmatrix}.
\]
The cross-ratio diameter and contraction factors are
\[
  \vartheta(A)=64,
  \qquad
  \kappa(A)=\frac79,
  \qquad
  \gamma=\frac{49}{81}.
\]
The source reports:
\[
\begin{array}{c|cc}
 k&\widehat\lambda_k&\lambda_k\\\hline
 0&2.000000&10.643722\\
 1&1.100000&1.418624\\
 2&1.015094&1.057195\\
 3&1.002393&1.008932\\
 4&1.000382&1.001424\\
 5&1.000061&1.000228
\end{array}
\]
Thus the a priori estimate is conservative initially but tracks the geometric
rate accurately after a few cycles.

\section*{Example 2: a much smaller contraction factor}
Take
\[
  p=q=\left(\frac13,\frac13,\frac13\right)^T,
  \qquad
  A=A_0=\frac1{30}
  \begin{pmatrix}
    3&4&4\\
    3&3&3\\
    4&3&4
  \end{pmatrix}.
\]
The limiting scaling is
\[
  \widehat B=
  \begin{pmatrix}
    0.093836321&0.125115095&0.114381917\\
    0.114381917&0.114381917&0.104569950\\
    0.125115095&0.093836321&0.114381917
  \end{pmatrix},
\]
and
\[
  \vartheta(A)=\frac{16}{9},
  \qquad
  \kappa(A)=\frac17,
  \qquad
  \gamma=\frac1{49}.
\]
The reported comparison is
\[
\begin{array}{c|cc}
 k&\widehat\lambda_k&\lambda_k\\\hline
 0&1.165685425&1.344914461\\
 1&1.001839973&1.002817612\\
 2&1.000001594&1.000002439\\
 3&1.000000001&1.000000002
\end{array}
\]
This example illustrates the practical force of the squared Birkhoff factor:
when $\kappa(A)$ is small, a full row--column cycle converges extremely rapidly.

\section*{Relation to projective-lag formulations}
The central theorem shape is not the entrywise bound by itself.
The load-bearing results are the vector-level contraction and the matrix increment
identity:
\[
  d(r^{(k+1)},p)\le \gamma d(r^{(k)},p),
  \qquad
  d(c^{(k+1)},q)\le \gamma d(c^{(k)},q),
\]
\[
  \Delta(A_k,A_{k+1})=d(r^{(k)},p)+d(c^{(k)},q).
\]
Together they imply asymptotic regularity of the diagonal correction vectors:
\[
  d\left(\frac{p}{r^{(k)}},e\right)\to 0,
  \qquad
  d\left(\frac{q}{c^{(k)}},e\right)\to 0.
\]
If a phase sequence is identified with one of these correction vectors,
then Hilbert-distance convergence to the constant ray gives the coordinate
cross-product collapse usually needed for projective lag:
\[
  u_{k,i}u_{k,j}^{-1}\to 1
  \quad\text{uniformly in }i,j,
\]
or, after clearing denominators in a bounded positive box,
\[
  u_{k,i}u_{k-1,j}-u_{k,j}u_{k-1,i}\to 0
\]
for the appropriate consecutive phase streams.

\paragraph{Hypotheses that cannot be dropped.}
The argument uses strict positivity of the kernel matrix to obtain a finite
multiplicative projective diameter and hence $\gamma<1$.  Bounds saying that the
iterates remain in a positive box are not enough by themselves; the contraction input
comes from the positive linear maps $A$ and $A^T$.

\paragraph{Two equivalent presentations.}
One can either follow Franklin--Lorenz close to the paper by defining
$A_k=S_kAT_k$ and proving the row/column normalization identities, or extract the
vector-level content of Lemma 2 and state it directly for the phase update maps.
The latter is closer to a phase-spread formulation; the former is closer to
the published theorem and gives a useful independent audit trail.

\begin{thebibliography}{99}
\bibitem{ChenGeorgiouPavon2021} Yongxin Chen, Tryphon T. Georgiou, and Michele Pavon. \emph{Stochastic Control Liaisons: Richard Sinkhorn Meets Gaspard Monge on a Schr\"odinger Bridge}. SIAM Review, 63(2):249--313, 2021. doi:10.1137/20M1339982. arXiv:2005.10963.
\bibitem{Leonard2014Survey} Christian L\'eonard. \emph{A Survey of the Schr\"odinger Problem and Some of Its Connections with Optimal Transport}. Discrete and Continuous Dynamical Systems--A, 34(4):1533--1574, 2014. doi:10.3934/dcds.2014.34.1533. arXiv:1308.0215.
\bibitem{Fortet1940} Robert Fortet. \emph{Resolution d'un systeme d'equations de M. Schrodinger}. Journal de mathematiques pures et appliquees, 9e serie, 19(1--4):83--105, 1940.
\bibitem{SinkhornKnopp1967} Richard Sinkhorn and Paul Knopp. \emph{Concerning Nonnegative Matrices and Doubly Stochastic Matrices}. Pacific Journal of Mathematics, 21(2):343--348, 1967.
\bibitem{FranklinLorenz1989} Joel Franklin and Jens Lorenz. \emph{On the Scaling of Multidimensional Matrices}. Linear Algebra and its Applications, 114/115:717--735, 1989. doi:10.1016/0024-3795(89)90490-4.
\bibitem{Carroll2004} Joseph E. Carroll. \emph{Birkhoff's Contraction Coefficient}. Linear Algebra and its Applications, 389:227--234, 2004. doi:10.1016/j.laa.2004.02.039.
\bibitem{EvesonNussbaum1995} Simon P. Eveson and Roger D. Nussbaum. \emph{An Elementary Proof of the Birkhoff--Hopf Theorem}. Mathematical Proceedings of the Cambridge Philosophical Society, 117(1):31--55, 1995.
\bibitem{BlanchetMurthy2019} Jose Blanchet and Karthyek R. A. Murthy. \emph{Quantifying Distributional Model Risk via Optimal Transport}. Mathematics of Operations Research, 44(2):565--600, 2019. doi:10.1287/moor.2018.0936. arXiv:1604.01446.
\bibitem{GaoKleywegt2023} Rui Gao and Anton J. Kleywegt. \emph{Distributionally Robust Stochastic Optimization with Wasserstein Distance}. Mathematics of Operations Research, 48(2):603--655, 2023. doi:10.1287/moor.2022.1275. arXiv:1604.02199.
\bibitem{MohajerinEsfahaniKuhn2018} Peyman Mohajerin Esfahani and Daniel Kuhn. \emph{Data-Driven Distributionally Robust Optimization Using the Wasserstein Metric: Performance Guarantees and Tractable Reformulations}. Mathematical Programming, 171(1--2):115--166, 2018. doi:10.1007/s10107-017-1172-1. arXiv:1505.05116.
\bibitem{WangGaoXie2025} Jie Wang, Rui Gao, and Yao Xie. \emph{Sinkhorn Distributionally Robust Optimization}. Operations Research, 2025. doi:10.1287/opre.2023.0294. arXiv:2109.11926.
\bibitem{Girsanov1960} I. V. Girsanov. \emph{On Transforming a Certain Class of Stochastic Processes by Absolutely Continuous Substitution of Measures}. Theory of Probability and Its Applications, 5(3):285--301, 1960. doi:10.1137/1105027.
\bibitem{CameronMartin1945} R. H. Cameron and W. T. Martin. \emph{Transformations of Wiener Integrals under a General Class of Linear Transformations}. Transactions of the American Mathematical Society, 58(2):184--219, 1945.
\bibitem{Yeh1978} J. Yeh. \emph{Transformation of Conditional Wiener Integrals under Translation and the Cameron--Martin Translation Theorem}. Tohoku Mathematical Journal, 30(4):505--515, 1978.
\bibitem{Kakutani1948} Shizuo Kakutani. \emph{On Equivalence of Infinite Product Measures}. Annals of Mathematics, 49(1):214--224, 1948.
\bibitem{ShiDeBortoliCampbellDoucet2023} Yuyang Shi, Valentin De Bortoli, Andrew Campbell, and Arnaud Doucet. \emph{Diffusion Schr\"odinger Bridge Matching}. Advances in Neural Information Processing Systems, 2023. arXiv:2303.16852.
\bibitem{LiuLipmanNickelKarrerTheodorouChen2024} Guan-Horng Liu, Yaron Lipman, Maximilian Nickel, Brian Karrer, Evangelos A. Theodorou, and Ricky T. Q. Chen. \emph{Generalized Schr\"odinger Bridge Matching}. International Conference on Learning Representations, 2024. arXiv:2310.02233.
\bibitem{DomingoEnrichDrozdzalKarrerChen2025} Carles Domingo-Enrich, Michal Drozdzal, Brian Karrer, and Ricky T. Q. Chen. \emph{Adjoint Matching: Fine-Tuning Flow and Diffusion Generative Models with Memoryless Stochastic Optimal Control}. arXiv preprint arXiv:2409.08861, 2025.

\bibitem{Birkhoff1957} Garrett Birkhoff. \emph{Extensions of Jentzsch's Theorem}. Transactions of the American Mathematical Society, 85(1):219--227, 1957.
\bibitem{Sinkhorn1964} Richard Sinkhorn. \emph{A Relationship Between Arbitrary Positive Matrices and Doubly Stochastic Matrices}. Annals of Mathematical Statistics, 35(2):876--879, 1964.
\bibitem{Sinkhorn1967Prescribed} Richard Sinkhorn. \emph{Diagonal Equivalence to Matrices with Prescribed Row and Column Sums}. American Mathematical Monthly, 74(4):402--405, 1967.
\bibitem{Bauer1965} Friedrich L. Bauer. \emph{An Elementary Proof of the Hopf Inequality for Positive Operators}. Numerische Mathematik, 7:331--337, 1965.
\bibitem{Bushell1973} P. J. Bushell. \emph{Hilbert's Metric and Positive Contraction Mappings in a Banach Space}. Archive for Rational Mechanics and Analysis, 52:330--338, 1973.

\end{thebibliography}

\end{document}
